\chapter{Methods}

\section{Study design}

We use a job choice survey experiment. Participants are repeatedly shown two hypothetical jobs,
side by side, and asked which they would prefer. Each job is described by a set of attributes ---
things like salary, working hours, and whether time off for appointments is easily accommodated
--- and the level of each attribute (e.g.\ ``40--44 hours per week'' vs.\ ``50--54 hours per
week'') is assigned at random each time a job is shown. Because the levels are randomly assigned,
we can isolate the causal effect of each individual attribute on whether a job is chosen, in the
same way a randomised trial isolates the effect of a treatment.

Neither job is labelled as a teaching job or a non-teaching job. Instead, attribute levels are
written in language that applies equally to teaching and other graduate occupations (for example,
``support staff to reduce workload'' rather than ``teaching assistant availability''). This lets
us model the decision to stay in or leave the profession, rather than only the decision to move
between schools.

We ask participants to imagine that two jobs are identical on everything except the attributes
shown, and simply state which they would prefer. This is a more demanding task than asking people
to rate arrangements on a scale, but it is also more informative: because salary is one of the
attributes tested, we can express how much of each other attribute is ``worth'' in salary terms
--- what we call its \emph{salary-equivalent value} throughout this report.

\section{What we asked about}

Table~\ref{table:attributes} sets out the family-friendly attributes we tested, alongside the
levels of each. Attributes were chosen to reflect arrangements that are common, or becoming more
common, in schools in England, which makes our findings directly actionable for schools looking
to expand flexible working. The exact wording was refined through cognitive interviews with
serving teachers and a pilot survey ahead of fieldwork, to check that each attribute was
understood as intended.

\begin{table}[htbp]
  \centering
  \caption{Family-friendly job attributes tested}
  \label{table:attributes}
  \begin{tabular}{p{4cm}p{9.5cm}}
    \toprule
    Attribute & Levels tested (low to high) \\
    \midrule
    Hours actually worked & 50--54 / 45--49 / 40--44 hours per week \\
    Working from home & None / half-day PPA from home / 2 days per week \\
    Part-time / job share & Discouraged / considered case-by-case / actively supported \\
    Compressed working week & Standard 5-day week / 9-day fortnight / 4-day week \\
    Time off for appointments & Discouraged / accommodated with advance planning / flexibly supported \\
    Work outside contracted hours & Regular evenings/weekends / occasional / rare \\
    Support staff for workload & None / some, for specific tasks / dedicated support regularly available \\
    Paid maternity leave & 66\% pay, 18 weeks / 100\%, 18 weeks / 100\%, 24 weeks \\
    Paid paternity leave & 90\%, 2 weeks / 100\%, 18 weeks / 100\%, 24 weeks \\
    Work--home fatigue spillover & Often interferes / sometimes interferes / rarely interferes \\
    Salary & 10\% lower / same / 10\% higher than current \\
    \bottomrule
  \end{tabular}
\end{table}

Because eleven attributes are too many to show at once without overwhelming respondents,
particularly on a mobile phone screen, each participant sees a subset in any one choice task
(a ``partial profile'' design). Salary appears in every task, which is what allows us to express
every other attribute in comparable salary-equivalent terms. Full detail on how attributes were
allocated across survey ``decks,'' and the robustness checks this requires, is reported in the
technical appendix.

\section{Who took part}

We administered the survey via Teacher Tapp, a daily panel survey completed by teachers across
England on their phones. We expect a sample of around 7,700 teachers, which gives us high
statistical power to detect the size of effect typically found in similar studies, and enough
teachers in their 30s specifically to support the subgroup comparisons at the heart of RQ3. Full
power calculations and a discussion of how representative the Teacher Tapp panel is of the wider
teacher workforce are set out in the technical appendix.

\section{How we analysed the data}

For each attribute, we estimate how much including a given level changes the probability that a
job is chosen, holding everything else constant --- what is known as the average marginal
component effect, or AMCE. We report these alongside marginal means, which show how appealing
each attribute level is on its own terms, without reference to a baseline category. Expressing
every attribute's AMCE relative to the AMCE of a 10\% salary increase gives us the
salary-equivalent value used throughout the findings. For RQ3, we compare these preferences
across subgroups (age, gender, parenting responsibilities, and stated intention to quit) using
conditional marginal means, following \citet{leeper2020measuring}. Full model specification and
robustness checks (deck order effects, profile order effects, and an attribute-composition test)
are reported in the technical appendix.

% NOTE (quick-and-dirty draft, [date]): citation keys (leeper2020measuring) need
% matching entries in references.bib. This section describes the design as planned
% in the submitted bid — update once piloting/cognitive interviews (Aug-Sep) have
% finalised exact attribute wording, and once actual achieved sample size is known.
